\section[Notebook 1983-10: traces of odd powers of $P - Q$]{Notebook 1983-10: the trace of the odd powers of a difference of two idempotents\protect\piste\protect\verifie{Notebook1983_10}}\label{sec:quillen-1983-10}

\emph{Context.} \carnet{1983}{10} works on index formulas in cyclic cohomology, the theory that Connes was developing at the time and to which Quillen would devote several papers. The index of a Toeplitz operator, or of a graded operator, is written there as the trace of a sufficiently large odd power of an operator built from an idempotent.

\emph{The page.} On page~31 (its page~154), Quillen observes that his two index formulas, the Toeplitz form and the graded form, reduce to the same type of expression: the trace of an odd power of a difference $e_1 - e_2$ of two idempotents. He proposes to study $\operatorname{tr}((e_1 - e_2)^{2n+1})$ and to try to show that it is an integer independent of $n$. The page poses the question; it does not answer it.

\emph{Intuition.} In infinite dimension, $e_1 - e_2$ is not trace class in general, and only a sufficiently large power is; the question is then deep. In finite dimension, all the traces exist, and the expected answer is $\operatorname{tr} e_1 - \operatorname{tr} e_2$, the difference of the ranks when working over a field of characteristic zero. The classical computation puts $A = P - Q$ and $B = P + Q - 1$: these two matrices anticommute and $A^2 + B^2 = 1$ (figure~\ref{fig:quillen-1983-10}). For $m$ odd, $\operatorname{tr}(A^m B^2)$ equals its opposite, by cyclicity of the trace and anticommutation, whence $\operatorname{tr}(A^m) = \operatorname{tr}(A^{m+2})$ \emph{provided one can divide by~$2$}. One might think that characteristic~$2$ is an exception. It is not.


\begin{figure}[htbp]
\centering
\begin{tikzpicture}[scale=1.05, >=stealth, every node/.style={font=\small}]
\draw[black!40, thick] (-2.6,0) -- (2.8,0) node[right] {$\ell_P$};
\draw[black!40, thick] ({-1.9*cos(50)},{-1.9*sin(50)}) -- ({2.9*cos(50)},{2.9*sin(50)}) node[above right=-2pt] {$\ell_Q$};
\draw[frise, dashed] ({-1.1*cos(25)},{-1.1*sin(25)}) -- ({2.8*cos(25)},{2.8*sin(25)}) node[right] {$b$};
\draw[piste, dashed] ({-2.4*cos(-20)},{-2.4*sin(-20)}) -- ({2.8*cos(-20)},{2.8*sin(-20)}) node[right] {$a$};
\draw (0.6,0) arc (0:50:0.6); \node at (13:0.85) {$\varphi$};
\draw[black!45, densely dotted] (-1.2,2.0785) -- (-1.2,0);
\draw[black!45, densely dotted] (-1.2,2.0785) -- (0.5276,0.6287);
\draw[->, thick] (0,0) -- (-1.2,2.0785) node[above] {$v$};
\draw[->, thick] (0,0) -- (-1.2,0) node[above left=-1pt] {$Pv$};
\draw[->, thick] (0,0) -- (0.5276,0.6287) node[above left=-1pt] {$Qv$};
\draw[->, thick, piste] (0,0) -- (-1.7276,-0.6287) node[below left=-2pt] {$Av$};
\draw[->, thick, frise] (0,0) -- (0.5276,-1.4498) node[below] {$Bv$};
\fill (0,0) circle (1.2pt);
\end{tikzpicture}
\caption{Two orthogonal projectors $P$ and $Q$ of the Euclidean plane, onto lines $\ell_P$ and $\ell_Q$ making an angle $\varphi$ (here $50^\circ$). The matrix $A = P - Q$ is $\sin\varphi$ times the orthogonal reflection in the axis $a$, and $B = P + Q - 1$ is $\cos\varphi$ times the reflection in the axis $b$, the bisector of $\ell_P$ and $\ell_Q$. The two axes make an angle of $45^\circ$: the product of the two reflections is a quarter turn, whence $AB = -BA$, and $A^2 + B^2 = (\sin^2\varphi + \cos^2\varphi)\cdot 1 = 1$. Here $\operatorname{tr} A^{2k+1} = 0 = \operatorname{tr} P - \operatorname{tr} Q$.}
\label{fig:quillen-1983-10}
\end{figure}

\begin{encadre}
\begin{theoreme}[the edition's statement; answer, in finite dimension, to the question of notebook 1983-10, p.~31]\label{thm:quillen-1983-10}
Let $R$ be a commutative ring and $P$, $Q$ two idempotent square matrices with entries in $R$. For every integer $k \ge 0$,
\[
\operatorname{tr}\bigl((P - Q)^{2k+1}\bigr) = \operatorname{tr} P - \operatorname{tr} Q.
\]
\end{theoreme}
\end{encadre}

\begin{proof}
Put $C = (P - Q)^2 = P + Q - PQ - QP$. Expanding, with $P^2 = P$ and $Q^2 = Q$,
\[
PC = P - PQP = CP, \qquad QC = Q - QPQ = CQ :
\]
$C$ commutes with $P$ and $Q$, hence so do its powers. It follows that
\[
(P - Q)\,C = (P - Q) - (PQP - QPQ).
\]
Let $D$ be a matrix commuting with $P$ and $Q$. By cyclicity of the trace, $\operatorname{tr}(DPQP) = \operatorname{tr}(PDPQ) = \operatorname{tr}(DP^2Q) = \operatorname{tr}(DPQ)$; likewise $\operatorname{tr}(DQPQ) = \operatorname{tr}(DQP)$, and $\operatorname{tr}(DQP) = \operatorname{tr}(PDQ) = \operatorname{tr}(DPQ)$. Hence
\[
\operatorname{tr}(DPQP) = \operatorname{tr}(DQPQ).
\]
Finally $(P - Q)^{2k+1} = C^k (P - Q)$, and
\[
(P - Q)^{2k+3} = C^k (P - Q)\, C = (P - Q)^{2k+1} - \bigl(C^k PQP - C^k QPQ\bigr).
\]
The trace of the term in parentheses is zero, by the above with $D = C^k$. The traces of the odd powers of $P - Q$ are therefore all equal to $\operatorname{tr}(P - Q) = \operatorname{tr} P - \operatorname{tr} Q$.
\end{proof}

\noindent The proof uses only idempotence, the commutation of $C$ with $P$ and $Q$, and the cyclicity of the trace; it never divides by~$2$. The computation by anticommutation shows $2\operatorname{tr}(A^m B^2) = 0$; the proof above shows directly $\operatorname{tr}(A^m B^2) = 0$, since $A^m B^2 = C^{k}(PQP - QPQ)$ for $m = 2k + 1$. An exhaustive computation over all pairs of idempotents of size~$2$ and~$3$ over $\mathbf F_2$, of size~$2$ over $\mathbf Z/4$, and over $\mathbf F_3$, $\mathbf F_5$, $\mathbf Z/9$, for $k = 1, 2, 3$, confirms the identity, in characteristic~$2$ as elsewhere.

\emph{Verification in Lean.} The file \href{\leanblobmain/Notebook1983_10.lean}{\texttt{Notebook1983\_10.\allowbreak{}lean}} proves Theorem~\ref{thm:quillen-1983-10} under the name \texttt{trace\_sub\_pow\_odd}, for matrices indexed by an arbitrary finite type over an arbitrary commutative ring, idempotence being mathlib's predicate \texttt{IsIdempotentElem}. The steps are \texttt{sq\_sub\_commute\_left} and \texttt{sq\_sub\_commute\_right} (commutation of $C$), \texttt{sub\_mul\_sq\_sub} (the identity $(P - Q)C = (P - Q) - (PQP - QPQ)$) and \texttt{trace\_mul\_PQP}. Mathlib does not contain this identity: a pattern search on \texttt{Matrix.trace} and \texttt{IsIdempotentElem} returns nothing. The axioms used are the three usual axioms of mathlib.

Not proved here: the infinite-dimensional version, where $P - Q$ is not trace class and only a sufficiently large odd power is, which is the real question of the page; the interpretation of the common value as the index of a Fredholm operator.

\litterature{the infinite-dimensional version is published: J.~Avron, R.~Seiler and B.~Simon \cite{ass1994} prove, for two orthogonal projections of a Hilbert space whose difference has a trace-class odd power $(P - Q)^{2n+1}$, that the index of the pair equals $\operatorname{tr}((P - Q)^{2n+1})$, and hence does not depend on $n$; their proof rests in particular on the anticommutation of $P - Q$ and $P + Q - 1$. The finite-dimensional case is its simplest example. The form with no hypothesis on the characteristic has not been found in print; it is elementary, and the edition does not present it as new. It is not a lead.}
